\documentclass[UTF8,11pt]{ctexart}
\usepackage[a4paper,margin=24mm]{geometry}
\usepackage{amsmath,amssymb,amsthm,mathtools,mathrsfs}
\usepackage[all]{xy}
\usepackage{xurl,hyperref,enumitem}
\usepackage{tikz}
\usetikzlibrary{arrows.meta,cd,decorations.markings,calc,patterns}
\hypersetup{hidelinks,breaklinks=true}
\setlength{\emergencystretch}{3em}
\allowdisplaybreaks
\newtheorem*{theorem}{Theorem}
\newtheorem*{thm}{Theorem}
\newtheorem*{lemma}{Lemma}
\newtheorem*{proposition}{Proposition}
\newtheorem*{prop}{Proposition}
\newtheorem*{corollary}{Corollary}
\newtheorem*{cor}{Corollary}
\newtheorem*{claim}{Claim}
\theoremstyle{definition}
\newtheorem*{definition}{Definition}
\newtheorem*{defn}{Definition}
\newtheorem*{remark}{Remark}
\newtheorem*{example}{Example}
\newcommand{\R}{\mathbb R}
\newcommand{\C}{\mathbb C}
\newcommand{\Z}{\mathbb Z}
\newcommand{\Q}{\mathbb Q}
\newcommand{\N}{\mathbb N}
\newcommand{\F}{\mathbb F}
\newcommand{\D}{\mathbb D}
\newcommand{\bB}{\mathbb B}
\newcommand{\bC}{\mathbb C}
\newcommand{\bR}{\mathbb R}
\newcommand{\bZ}{\mathbb Z}
\newcommand{\bN}{\mathbb N}
\newcommand{\bQ}{\mathbb Q}
\newcommand{\bD}{\mathbb D}
\newcommand{\bF}{\mathbb F}
\newcommand{\CP}{\mathbb{CP}}
\newcommand{\RP}{\mathbb{RP}}
\newcommand{\sO}{\mathscr O}
\newcommand{\sI}{\mathscr I}
\newcommand{\sF}{\mathscr F}
\newcommand{\sG}{\mathscr G}
\newcommand{\sH}{\mathscr H}
\newcommand{\sR}{\mathscr R}
\newcommand{\sM}{\mathscr M}
\newcommand{\sV}{\mathscr V}
\newcommand{\sU}{\mathscr U}
\DeclareMathOperator{\Hom}{Hom}
\DeclareMathOperator{\Ext}{Ext}
\DeclareMathOperator{\Tor}{Tor}
\DeclareMathOperator{\Spec}{Spec}
\DeclareMathOperator{\Proj}{Proj}
\DeclareMathOperator{\supp}{supp}
\DeclareMathOperator{\im}{im}
\DeclareMathOperator{\id}{id}
\DeclareMathOperator{\Tr}{Tr}
\DeclareMathOperator{\tr}{tr}
\DeclareMathOperator{\rank}{rank}
\DeclareMathOperator{\ord}{ord}
\DeclareMathOperator{\dist}{dist}
\DeclareMathOperator{\End}{End}
\DeclareMathOperator{\Aut}{Aut}
\DeclareMathOperator{\Gal}{Gal}
\DeclareMathOperator{\vol}{vol}
\DeclareMathOperator{\Cl}{Cl}
\DeclareMathOperator{\coker}{coker}
\DeclareMathOperator{\codim}{codim}
\DeclareMathOperator{\divergence}{div}
\DeclareMathOperator{\Ric}{Ric}
\DeclareMathOperator{\grad}{grad}
\DeclareMathOperator{\Hess}{Hess}
\newcommand{\abs}[1]{\left\lvert#1\right\rvert}
\newcommand{\sabs}[1]{\lvert#1\rvert}
\newcommand{\babs}[1]{\bigl\lvert#1\bigr\rvert}
\newcommand{\Babs}[1]{\Bigl\lvert#1\Bigr\rvert}
\newcommand{\bbabs}[1]{\biggl\lvert#1\biggr\rvert}
\newcommand{\BBabs}[1]{\Biggl\lvert#1\Biggr\rvert}
\newcommand{\norm}[1]{\left\lVert#1\right\rVert}
\newcommand{\snorm}[1]{\lVert#1\rVert}
\newcommand{\bnorm}[1]{\bigl\lVert#1\bigr\rVert}
\newcommand{\Bnorm}[1]{\Bigl\lVert#1\Bigr\rVert}
\newcommand{\bbnorm}[1]{\biggl\lVert#1\biggr\rVert}
\newcommand{\BBnorm}[1]{\Biggl\lVert#1\Biggr\rVert}
\newcommand{\widebar}[1]{\overline{#1}}
\newcommand{\nicefrac}[2]{\frac{#1}{#2}}
\newcommand{\linnprod}[2]{\langle#1,#2\rangle}
\newcommand{\myindex}[1]{#1}
\newcommand{\glsadd}[1]{}
\newcommand{\avoidbreak}{}
\newcommand{\sourceref}[1]{\textnormal{[原文：\nolinkurl{#1}]}}
\newcommand{\thmref}[1]{\sourceref{#1}}
\newcommand{\lemmaref}[1]{\sourceref{#1}}
\newcommand{\propref}[1]{\sourceref{#1}}
\newcommand{\corref}[1]{\sourceref{#1}}
\newcommand{\defnref}[1]{\sourceref{#1}}
\newcommand{\exampleref}[1]{\sourceref{#1}}
\newcommand{\exerciseref}[1]{\sourceref{#1}}
\newcommand{\figureref}[1]{\sourceref{#1}}
\newcommand{\sectionref}[1]{\sourceref{#1}}
\newcommand{\appendixref}[1]{\sourceref{#1}}
\newcommand{\stacksref}[2]{\href{https://stacks.math.columbia.edu/tag/#1}{#2 [#1]}}

\begin{document}
\section*{041\quad 辛形式同痕的Moser构造与Darboux局部标准形}
\subsection*{来源与计数目标}
Kartik Venkatram 与 Denis Auroux 合作整理的 MIT 18.966 \emph{Geometry of Manifolds} 人类课程讲义，Spring 2007。课程页面明确记载该署名关系；不是将学生笔记冒充授课者独著。Lecture 4，印刷第2页，Definition 2、Theorem 1 与 Theorem 2 及完整证明。获取日期：2026-09-28。
\par\url{https://ocw.mit.edu/courses/18-966-geometry-of-manifolds-spring-2007/864a8cbd55d8d9696a717db38f8668a9_lect04.pdf}
\par\url{https://ocw.mit.edu/courses/18-966-geometry-of-manifolds-spring-2007/pages/lecture-notes/}
\par\textbf{本文件只计一个问题：}构造保持辛形式族的同痕，并以同一构造取得 Darboux 坐标；合并计一个问题。
\subsection*{作者命题与人类证明}

\paragraph{2. Moser's Theorem.}
One can ask whether, for a given manifold $M$, two symplectic structures $\omega_0,\omega_1$ are equivalent, i.e. whether there is a symplectomorphism $M\to M$ which pulls back one to the other. In general, $[\omega_0]=[\omega_1]$ does not imply that the two structures are symplectomorphic. To study this question further, we give other notions of equivalence.
\begin{definition}[2]
Two forms $\omega_0,\omega_1$ are deformation equivalent if $\exists(\omega_t)_{t\in[0,1]}$ a continuous family of symplectic forms, and isotopic if there is such a family with $[\omega_t]$ constant in $H^2(M,\R)$.
\end{definition}
Let $M$ be a compact manifold with $\omega_0,\omega_1$ isotopic symplectic forms (i.e. $\exists\omega_t$ as above with each $\omega_t$ nondegenerate).
\begin{theorem}[1: Moser]
$\exists$ an isotopy $\rho_t:M\to M$ s.t. $\rho_t^*\omega_t=\omega_0$.
\end{theorem}
That is, $(M,\omega_0)$ and $(M,\omega_1)$ are symplectomorphic.
\begin{proof}
(This technique is known as Moser's trick.) By assumption, $[\omega_t]$ is independent of $t$, i.e. $[\frac{d\omega_t}{dt}]=0$. Thus, $\exists\alpha_t$ a 1-form s.t. $\frac{d\omega_t}{dt}=-d\alpha_t$: moreover, we can choose this $\alpha_t$ smoothly w.r.t. to $t$ (via the Poincaré lemma). Since $\omega_t$ is nondegenerate, $\exists X_t$ s.t. $i_{X_t}\omega_t=\alpha_t$. Moreover, since $M$ is compact, we have a well-defined flow $\rho_t$ of $X_t$. Now,
\begin{align*}
\frac{d}{dt}(\rho_t^*\omega_t)
&=\rho_t^*(L_{X_t}\omega_t)+\rho_t^*\left(\frac{d\omega_t}{dt}\right)\\
&=\rho_t^*\left(di_{X_t}\omega_t+\frac{d\omega_t}{dt}\right)=0.
\end{align*}
Since $\rho_0$ is the identity, we have our desired isotopy.
\end{proof}

\begin{theorem}[2: Darboux]
For $(M,\omega)$ symplectic, $p\in M$, $\exists U\ni p$ with a coordinate system $(x_1,y_1,\ldots,x_n,y_n)$ s.t. $\omega|_U=\sum dx_i\wedge dy_i$.
\end{theorem}
\begin{proof}
$(T_pM,\omega_p)$ has a standard basis $(e_1,\ldots,e_n,f_1,\ldots,f_n)$, so there exist local coordinates $(x_1,y_1,\ldots,x_n,y_n)$ s.t. $\omega_p=\sum dx_i\wedge dy_i$. On a neighborhood $U$ of $p$, we obtain two symplectic forms: $\omega$ and the standard form. The family $\omega_t=(1-t)\omega_0+t\omega_1$ is one of closed forms: since nondegeneracy is an open condition, we can shrink our neighborhood to assure that $\omega_t$ is nondegenerate for each $t$ on some $U'\ni p$. Thus, $\exists\alpha\in\Omega^1(U)$ s.t. $\omega_1-\omega_0=-d\alpha$. Subtracting a constant, we can assume $\alpha_p=0$. Let $v_t$ be the vector field on $U$ s.t. $i_{v_t}\omega_t=\alpha$. Then $\exists U''\ni p$ s.t. its flow $\rho_t$ is defined $\forall t$. By the Moser's trick, we fnd that $\rho_1^*\omega_1=\omega_0$, implying that the symplectic form is indeed standard after composing our chosen coordinates with $\rho_1$.
\end{proof}

\subsection*{整理说明（非作者正文）}
按作者证明实际使用的情形，形式族随参数光滑；原定义用了 ``continuous family''，本整理保留该原句并在此明示证明使用的光滑参数范围。允许先修为 de Rham 上同调、带参数原始形式的选择、Cartan 公式、时变向量场的局部流与紧流形上的有限时间存在性、辛向量空间标准基。核心是从上同调不变性选取一形式，再以非退化性解出时变向量场并积分为所需坐标变换，不是直接套用 Darboux 或 Moser 定理。两个作者定理合并一次计数，不将局部特例另算。原第2页图像已读取；按原文字与公式转录，改写断行，保留 ``fnd'' 原拼写。
\subsection*{可修改的简短标注（非作者正文）}
$c$：辛形式的同痕等价与局部标准形。

$a$：de Rham 上同调；原始一形式；辛形式的非退化配对；时变向量场；Cartan 公式；流的拉回微分。

\texttt{cross\_theory\_status = yes}。上同调中的差分条件经收缩方程变为向量场，再经常微分方程的流产生保持形式的同痕；局部构造把这一过程返回坐标标准化。位置：Lecture 4 第2页两个证明。

全部标注均为非穷尽、可修改初稿；未标注不表示未使用。
\subsection*{许可与修改记录}
材料来自 MIT OpenCourseWare，作者与课程见来源。按该站所示 CC BY-NC-SA 4.0 许可复用，整理部分沿用同一许可。\url{https://creativecommons.org/licenses/by-nc-sa/4.0/}。2026-09-28：助手转录题证、调整数学排版并加入来源、范围和初稿标注；不暗示作者或 MIT 认可本次整理。所留为本次题证转录摘录，不是原 PDF 下载件。
\end{document}
